\setcounter{section}{0} % tema número N-1
\section{Almacenamiento}

Nombre completo: Almacenamiento para sistemas grandes y departamentales. Dispositivos para trataimento de información multimedia. Virtualización del almacenamiento. Copias de seguridad.

\localtableofcontents

\subsection{Soportes o Medios de almacenamiento de datos}
    \subsubsection{Discos}
        \paragraph{Discos magnéticos}
            Dispositivos basados en un estrato magnético giratorio. Incorporan una controladora que incluye una caché. Tienen gran capacidad (del orden de TB) y un tiempo de acceso medio. Las últimas interfaces utilizadas son \textbf{SATA} 600MB/s y \textbf{SAS} 2400MB/s

        \paragraph{Discos de estado sólido}
            Basan su almacenamiento en memoria no volátil. Tienen mayor ratio de lectura y escritura, menos peso y tamaño, a cambio son mas caros y tienen menos capacidad. Pueden usar las interfaces de los discos duros magnéticos o \textbf{NVMe} y \textbf{PCI-e}, pensadas para reducir la latencia y aumentar el ancho de banda.

        \paragraph{Cabinas de discos}
            Dispositivos que aglutinan discos. Proporcionan rendimiento mediante caché e integran múltiples discos en una unidad lógica mediante RAID.

            En el fron-end permite la conexión con servidores por bloques, o ficheros. Se caracterizan por número de puertos y velocidad.

            En el back-end se ofrece la conexion a los fistintos discos, típicamente FC y SAS.

            En medio se encuentra la controladora que gestiona el funcionamiento de la cabina y proporciona características como RAID, cifrado, deduplicación, snapshots, backups, monitorización\dots

            Los volúmenes se pueden crear \textbf{thick-provisioning} donde se asigna el espacio de uso, o \textbf{thin-provisioning} donde se asignan dinámicamente los bloques que se vayan consumiento. Los datos pueden moverse con \textbf{autotiering} donde se acercan los datos mas frecuentes y se alejan los mas raros a discos mas rápidos o lentos. Por último se usa \textbf{deduplicación} que permite reducir los datos ya existentes, ideal para backups y repositorios.
    \subsubsection{RAID}
        Tecnologías basadas en la combinación de varios discos duros para formar una única unidad lógica. Se utilizan diferentes técicas y mecanismos.

        Se usa el \textbf{particionado} donde se divide la escritura en distintos discos, aumentando la velocidad y la tasa de transferencia.

        Se \textbf{duplican} datos para permitir integridad y tolerancia a fallos. Se usa \textbf{mirroring} para almacenarlos simultáneamente en varios discos, o \textbf{paridad} para obtener función resumen, que permite reconstruir los datos en caso de que un disco se pierda, esta necesita menos espacio que el mirroring.

        \begin{table}[H]
        \begin{tabular}{rlll}
            nivel & descripción & discos mínimos & Pérdida soportada\\
            hline
            RAID 0 & Stripping a nivel de bloque, pegar con cinta dos discos & 2 & 0 \\
            RAID 1 & Mirroring a nivel de bloque, un copiapega. La capacidad es la del disco mas pequeño.  & 2 & 1 \\
            RAID 2 & Tripping a nivel de bit con un disco de paridad. Todos los discos giran a la misma velocidad, en desuso & 3 & 1 \\
            RAID 3 & Stripping a nivel de byte con disco de paridad dedicado & 3 & 1 \\
            RAID 4 & Stripping a nivel de bloque cdon disco de paridad dedicado & 3 & 1\\
            RAID 5 & Stripping a nivel de bloque con paridad distribuida. La eficiencia es 1-1/n & 3 & 1 \\
            RAID 6 & STripping a nivel de bloque con paridad distribuida, (2 bloques) & 4 & 2\\
            RAID 0+1 & Espejo de dos RAID 0 & 4 & 2* \\
            RAID 1+0 & Stripping de dos mirror & 4 & 2*\\
        \end{tabular}
        \end{table}
        Las RAID 1+0 y 0+1 soportan hasta 2 díscos de pérdida, según cuales. Otras opciones de RAID como suma utilizan otras RAID como discos.

        También existen RAID 5E y 6E con discos de reserva como hot-spare que permiten minimizar el tiempo de reconstrucción en caso de fallo.



\subsection{Sistemas de almacenamiento de datos}

\subsection{Virtualización del almacenamiento y backup}

\subsection{Nuevas tendencias}